% --- Chapter: Time and Date Functions ---
\section{Time and Date Functions}

Working with time is essential in scripting: timestamps for logging, date formatting for reports, scheduling, and measuring elapsed time. CRISP provides six time functions that cover Unix timestamps, human-readable dates, formatting, and parsing—all without external dependencies.

\subsection{Getting the Current Time}

Three functions return the current time at different resolutions:

\begin{table}[H]
\centering
\rowcolors{2}{softviolet!30}{white}
\caption{\textcolor{darkpurple}{Time Functions}}
\vspace{0.3em}
\begin{tabular}{@{}lllp{6cm}@{}}
\toprule
\rowcolor{softvioletdark}
\textbf{\textcolor{white}{Function}} & \textbf{\textcolor{white}{Returns}} & \textbf{\textcolor{white}{Resolution}} & \textbf{\textcolor{white}{Example Output}} \\
\midrule
\texttt{time()} & Int & Seconds & \texttt{1784987128} \\
\texttt{timestamp()} & Int & Milliseconds & \texttt{1784987128638} \\
\texttt{datetime()} & String & Seconds & \texttt{"2026-07-25 15:45:28"} \\
\texttt{datetime\_utc()} & String & Seconds & \texttt{"2026-07-25 13:45:28 UTC"} \\
\bottomrule
\end{tabular}
\end{table}

\vspace{1.5cm}
\begin{tcolorbox}[colback=codebg, colframe=darkpurple, colbacktitle=darkpurple, title=\textbf{\textcolor{codegold}{Basic Time Functions}}, fonttitle=\bfseries, sharp corners]
\begin{lstlisting}
# Unix timestamp — seconds since epoch (Jan 1, 1970)
let now = time();
say "Seconds since epoch: " + now;

# Millisecond timestamp — for high-resolution timing
let ms = timestamp();
say "Milliseconds since epoch: " + ms;

# Human-readable local time
say "Local: " + datetime();
# "Local: 2026-07-25 15:45:28"

# Human-readable UTC time
say "UTC:   " + datetime_utc();
# "UTC:   2026-07-25 13:45:28 UTC"
\end{lstlisting}
\end{tcolorbox}

\clearpage
\subsection{Formatting Dates}

\texttt{strftime} formats a Unix timestamp into a human-readable string using format specifiers. This is the same \texttt{strftime} familiar to C, Perl, and Python programmers:

\begin{tcolorbox}[colback=codebg, colframe=darkpurple, colbacktitle=darkpurple, title=\textbf{\textcolor{codegold}{Date Formatting with strftime}}, fonttitle=\bfseries, sharp corners]
\begin{lstlisting}
let now = time();

# Common formats
say strftime("%Y-%m-%d", now);           # 2026-07-25  (ISO date)
say strftime("%H:%M:%S", now);           # 15:45:28    (time)
say strftime("%Y-%m-%d %H:%M:%S", now);  # 2026-07-25 15:45:28 (datetime)
say strftime("%A, %B %d, %Y", now);      # Saturday, July 25, 2026
say strftime("%I:%M %p", now);           # 03:45 PM    (12-hour)

# Custom log format
say strftime("[%Y-%m-%d %H:%M:%S]", now) + " Server started";
# [2026-07-25 15:45:28] Server started
\end{lstlisting}
\end{tcolorbox}

\vspace{1.5cm}
\subsection{Parsing Dates}

\texttt{strptime} does the reverse: it parses a date string into a Unix timestamp using the same format specifiers:

\begin{tcolorbox}[colback=codebg, colframe=darkpurple, colbacktitle=darkpurple, title=\textbf{\textcolor{codegold}{Date Parsing with strptime}}, fonttitle=\bfseries, sharp corners]
\begin{lstlisting}
# Parse a date string into a timestamp
let ts = strptime("2026-01-15", "%Y-%m-%d");
say ts;  # 1768435200 (Unix timestamp for Jan 15, 2026)

# Verify by formatting back
say strftime("%A, %B %d, %Y", ts);
# Thursday, January 15, 2026

# Parse a full datetime
let ts2 = strptime("2026-07-25 15:45:28", "%Y-%m-%d %H:%M:%S");
say ts2;  # should match time()
\end{lstlisting}
\end{tcolorbox}

\clearpage
\subsection{A Complete Example: Elapsed Time Calculator}

Measuring how long something takes is a classic scripting task. Here is a program that times a computation and reports the result in human-readable form:

\begin{tcolorbox}[colback=codebg, colframe=darkpurple, colbacktitle=darkpurple, title=\textbf{\textcolor{codegold}{Elapsed Time Calculator (timer.csp)}}, fonttitle=\bfseries, sharp corners]
\begin{lstlisting}
# timer.csp — Measure and display elapsed time

fn elapsed(start_ms, end_ms) {
    let diff_ms = end_ms - start_ms;

    if diff_ms < 1000 {
        return diff_ms + " ms";
    }

    let seconds = diff_ms / 1000.0;

    if seconds < 60 {
        return seconds + " seconds";
    }

    let minutes = int(seconds / 60.0);
    let secs_left = seconds - (minutes * 60.0);
    return minutes + " min " + secs_left + " sec";
}

# Time a computation
say "Computing factorial(20)...";
let start = timestamp();

let result = 1;
let i = 1;
while i <= 20 {
    result = result * i;
    i = i + 1;
}

let end = timestamp();

say "factorial(20) = " + result;
say "Took: " + elapsed(start, end);
\end{lstlisting}
\end{tcolorbox}

\clearpage
\subsection{Date Arithmetic}

Here is a practical program that computes age from a birth date and determines the day of the week someone was born:

\begin{tcolorbox}[colback=codebg, colframe=darkpurple, colbacktitle=darkpurple, title=\textbf{\textcolor{codegold}{Date Arithmetic (age\_calculator.csp)}}, fonttitle=\bfseries, sharp corners]
\begin{lstlisting}
# age_calculator.csp — Compute age and birth day-of-week

fn calculate_age(birth_date_str) {
    let birth_ts = strptime(birth_date_str, "%Y-%m-%d");
    let now_ts = time();

    # Seconds in a year (approximate)
    let seconds_per_year = 365.25 * 24 * 60 * 60;
    let age = int((now_ts - birth_ts) / seconds_per_year);

    let day_of_week = strftime("%A", birth_ts);

    return {
        age => age,
        day_of_week => day_of_week,
        birth_date => strftime("%B %d, %Y", birth_ts)
    };
}

let info = calculate_age("1990-05-15");
say "Born: " + info.birth_date;
say "Day of week: " + info.day_of_week;
say "Age: " + info.age + " years";
\end{lstlisting}
\end{tcolorbox}

\begin{tcolorbox}[colback=lightlavender, colframe=softvioletdark, title=\textbf{\textcolor{darkpurple}{Design Note: Why strftime/strptime?}}, fonttitle=\bfseries, sharp corners]
CRISP delegates date formatting and parsing to the C standard library's \texttt{strftime} and \texttt{strptime} functions. This is a deliberate choice:

\begin{itemize}
    \item \textbf{Familiarity:} The format specifiers (\texttt{\%Y}, \texttt{\%m}, \texttt{\%d}, \texttt{\%H}, \texttt{\%M}, \texttt{\%S}) are identical to C, Perl, Python, Ruby, and PHP. Anyone who has formatted dates in any of these languages already knows CRISP's date formatting.
    \item \textbf{Locale-aware:} \texttt{\%A} produces the correct weekday name for the system's locale. On a Czech system, \texttt{strftime("\%A", ts)} returns ``pondělí'', not ``Monday''.
    \item \textbf{No dependencies:} Using the C standard library means zero additional crates, zero build time, and zero security vulnerabilities from third-party date libraries.
\end{itemize}

The trade-off is that \texttt{strptime} is not available on Windows (it's a POSIX extension). For cross-platform date parsing, future CRISP releases will provide a pure-Rust alternative.
\end{tcolorbox}
